\begin{table}[tbp]
\centering
\caption{问题一：特征组消融}
\label{tab:t04_q1_ablation}
\footnotesize
\setlength{\tabcolsep}{3.5pt}
\begin{tabular}{lllrrrrl}
\toprule
特征组 & 处理口径 & 划分 & 平均 R² & 种子 SD & Δ vs 完整 & 95\% CI & 判定 \\
\midrule
完整（哨兵屏蔽 + 检出率协变量） & qc & 随机 CV & 0.2476 & 0.0019 & — & — & 有贡献或中性 \\
完整（哨兵屏蔽 + 检出率协变量） & qc & 分组 CV & 0.1785 & 0.0028 & — & — & 有贡献或中性 \\
去掉检出率协变量 & qc & 随机 CV & 0.2470 & 0.0023 & -0.0006 & [-0.002, 0.001] & 有贡献或中性 \\
去掉检出率协变量 & qc & 分组 CV & 0.1797 & 0.0037 & +0.0011 & [0.000, 0.002] & 有贡献或中性 \\
去掉派生汇总列 & qc & 随机 CV & 0.2451 & 0.0016 & -0.0025 & [-0.005, -0.000] & 有贡献或中性 \\
去掉派生汇总列 & qc & 分组 CV & 0.1758 & 0.0005 & -0.0027 & [-0.005, -0.000] & 有贡献或中性 \\
哨兵当真实值（不屏蔽） & plain & 随机 CV & 0.2485 & 0.0022 & +0.0009 & [-0.000, 0.002] & 有贡献或中性 \\
哨兵当真实值（不屏蔽） & plain & 分组 CV & 0.1797 & 0.0017 & +0.0012 & [-0.000, 0.002] & 有贡献或中性 \\
审计前口径（哨兵当真实值、无检出率） & plain & 随机 CV & 0.2491 & 0.0027 & +0.0015 & [0.000, 0.003] & 对照基线（仅参考） \\
审计前口径（哨兵当真实值、无检出率） & plain & 分组 CV & 0.1811 & 0.0013 & +0.0026 & [0.001, 0.004] & 对照基线（仅参考） \\
\bottomrule
\end{tabular}
\end{table}
% 注（写入图注/表注，勿直接编译）：Δ = 去掉该特征组后的平均 OOF R² − 完整模型；CI 为 3 个种子 × 50 维的 Bootstrap 区间。“有贡献或中性”表示 CI 未显示明确损失，不等同于该组必要。
